% CPSY 1291 -- Recitation handout 5
% Compile:  latexmk -lualatex handout-05-pytorch-training-loop.tex
\documentclass[11pt]{article}
\usepackage{cpsy1291-handout}

\begin{document}

\handouthead{5}{PyTorch: tensors, autograd, and the training loop}
  {The five lines that train every network in this course}
  {Week 5 \ $\cdot$\ hold this after Lecture 8; Assignment 2 is due Tue 10/13}

\section*{What this session is for}
\addcontentsline{toc}{section}{What this session is for}

Last week you derived a learning rule by hand and checked it numerically. This
week you hand that job to a library.

You will not be asked to write a training loop from scratch in this course. You
\emph{will} be asked to read one, modify one, and diagnose one that is not
working --- and to know which of its lines are doing something and which are
ceremony. That is the whole of this session.

\begin{keybox}
\ask{The one idea.} PyTorch is NumPy that \emph{remembers what you did to it}. Every
operation on a tensor is recorded, so when you finally produce one number --- the
loss --- the library can replay the recording backwards and hand you the
derivative with respect to every parameter that took part. That replay is
\texttt{loss.backward()}, and it is the only thing PyTorch does that NumPy
cannot.
\end{keybox}

\section{Tensors}

A tensor is an array. Almost every NumPy habit transfers: indexing, slicing,
\texttt{.shape}, broadcasting, \texttt{axis=} (spelled \texttt{dim=}).

\begin{lstlisting}
import torch

x = torch.tensor([[1., 2.], [3., 4.]])   # from a Python list
x = torch.from_numpy(arr).float()        # from a NumPy array -- note .float()
x.numpy()                                # and back again
x.shape                                  # torch.Size([2, 2])
x.mean(dim=0)                            # 'axis' is spelled 'dim'
\end{lstlisting}

\subsection{The dtype trap}

NumPy defaults to \texttt{float64}. PyTorch defaults to \texttt{float32} and
most layers refuse anything else. So converted arrays need an explicit cast:

\begin{lstlisting}
X = torch.tensor(X_np, dtype=torch.float32)      # do this
\end{lstlisting}

\begin{keybox}
\ask{The error you will get, and what it really means.}
\texttt{RuntimeError: expected scalar type Float but found Double}. ``Double''
is \texttt{float64}, i.e.\ a NumPy array that came across without a cast. It is
never a bug in your model.
\end{keybox}

A second dtype rule: classification \emph{targets} must be \texttt{long}
(integers), not floats, and they are class indices --- $0,1,2$ --- not one-hot
vectors.

\subsection{Batch first, always}

Every layer in PyTorch expects the first dimension to be the batch. A single
example still needs that dimension:

\begin{lstlisting}
x.shape             # (784,)      one image -- layers will reject this
x[None].shape       # (1, 784)    now it is a batch of one
\end{lstlisting}

The corresponding shapes for the data you will meet this term:

\begin{center}
\begin{tabular}{@{}lll@{}}
\toprule
Data & Shape & Read as \\
\midrule
tabular / flattened & $(N, D)$ & $N$ examples, $D$ features \\
images & $(N, C, H, W)$ & $N$ images, $C$ channels, height, width \\
sequences & $(N, T, D)$ & $N$ sequences, $T$ steps, $D$ features \\
\bottomrule
\end{tabular}
\end{center}

Note that images are \textbf{channels-first} in PyTorch and channels-\emph{last}
in matplotlib. \texttt{plt.imshow} on a $(3, 64, 64)$ tensor fails; you need
\texttt{img.permute(1, 2, 0)}.

\section{Autograd in three lines}

\begin{lstlisting}
w = torch.tensor([2.0], requires_grad=True)
loss = (w ** 2).sum()      # a graph is being recorded
loss.backward()            # replay it backwards
w.grad                     # tensor([4.])   -- d(w^2)/dw = 2w = 4
\end{lstlisting}

Three facts that explain most autograd confusion:

\begin{enumerate}
  \item \textbf{Gradients accumulate.} \texttt{.grad} is \emph{added} to, not
        replaced, on every \texttt{backward()}. This is deliberate --- it lets you
        sum gradients over several batches --- but it means you must call
        \texttt{opt.zero\_grad()} each step or you are descending on the sum of
        every gradient since the program started. Forget it and the loss
        typically diverges after a few epochs.
  \item \textbf{\texttt{backward()} only works from a scalar.} The loss must be
        one number. If it is not, you forgot a \texttt{.mean()} or a
        \texttt{.sum()}.
  \item \textbf{The graph is freed after the backward pass.} Calling
        \texttt{backward()} twice on the same loss raises an error. That error
        almost always means you built the loss outside the loop.
\end{enumerate}

\begin{keybox}
\ask{\texttt{torch.no\_grad()}.} Inside this block nothing is recorded. Use it
whenever you are evaluating rather than training --- computing validation
accuracy, extracting features from a pretrained network, plotting. It is faster
and it uses far less memory, and forgetting it is the usual reason a script runs
out of memory during evaluation but not during training.
\end{keybox}

\section{Building a model}

\subsection{\texttt{nn.Linear} is the neuron layer from lecture}

\texttt{nn.Linear(in\_features, out\_features)} computes
$\boldsymbol{x}\boldsymbol{W}^{\!\top} + \boldsymbol{b}$ --- a whole layer of the
units you studied last week, in one object. Its \texttt{weight} attribute has
shape \texttt{(out, in)}, which is the transpose of what you might expect; that
is a convention, and reading it wrong is a common source of confusion when you
inspect a trained model.

\begin{lstlisting}
import torch.nn as nn

model = nn.Sequential(
    nn.Linear(64, 16),      # 64 inputs -> 16 hidden units
    nn.ReLU(),
    nn.Linear(16, 10),      # 16 -> 10 class scores
)
sum(p.numel() for p in model.parameters())     # count the parameters
\end{lstlisting}

\texttt{nn.Sequential} is enough for everything in Assignment 2. Note that
\texttt{ReLU} is a layer here, sitting \emph{between} the linear layers ---
without it the whole stack collapses to a single linear map, which is exactly
the point Lecture 8 makes about why depth needs nonlinearity.

\subsection{The output layer, and the trap}

The last layer emits \textbf{logits} --- raw, unbounded scores, one per class.
They are not probabilities and they are not meant to be.

\begin{keybox}
\ask{Never put a ReLU or a softmax after your final layer if you are using
\texttt{nn.CrossEntropyLoss}.} That loss applies \texttt{log\_softmax}
internally. Applying softmax yourself first applies it twice and the model
trains badly for reasons nothing will report. Applying \emph{ReLU} first is
worse: it clamps every negative logit to zero, so the model can no longer
express ``this class is unlikely'' at all. Both are silent. Both are common in
code found online.
\end{keybox}

\begin{center}
\begin{tabular}{@{}lll@{}}
\toprule
Task & Loss & Final layer emits \\
\midrule
regression & \texttt{nn.MSELoss()} & one number per target, no activation \\
classification & \texttt{nn.CrossEntropyLoss()} & raw logits, no activation \\
binary, if you must & \texttt{nn.BCEWithLogitsLoss()} & one raw logit \\
\bottomrule
\end{tabular}
\end{center}

\section{The training loop}

Here it is, complete. Read it as five jobs, not as eleven lines.

\begin{lstlisting}
model = nn.Sequential(nn.Linear(64, 16), nn.ReLU(), nn.Linear(16, 10))
lossfn = nn.CrossEntropyLoss()
opt    = torch.optim.Adam(model.parameters(), lr=1e-3)

for epoch in range(200):
    model.train()
    opt.zero_grad()               # 1. forget last step's gradients
    out  = model(Xtr)             # 2. forward: predictions
    loss = lossfn(out, ytr)       # 3. how wrong are we?
    loss.backward()               # 4. gradient of loss w.r.t. every parameter
    opt.step()                    # 5. take one step downhill

    if epoch % 20 == 0:
        model.eval()
        with torch.no_grad():
            acc = (model(Xte).argmax(1) == yte).float().mean()
        print(f"{epoch:4d}  loss {loss.item():.4f}  test acc {acc:.3f}")
\end{lstlisting}

\begin{itemize}
  \item \texttt{opt.zero\_grad()} --- because gradients accumulate.
  \item \texttt{loss.backward()} --- the replay from the box on page 1.
  \item \texttt{opt.step()} --- \emph{this} is the line that changes the weights.
        Steps 1--4 only compute; nothing moves until here.
  \item \texttt{.item()} --- pulls a Python float out of a one-element tensor.
        Printing the tensor itself works but drags the whole graph along.
  \item \texttt{model.train()} / \texttt{model.eval()} --- switches the behavior
        of dropout and batch-norm layers. Harmless with the model above; a real
        bug the moment you use either, and free to write, so write it.
\end{itemize}

\subsection{SGD or Adam?}

\begin{center}
\begin{tabular}{@{}lll@{}}
\toprule
& \texttt{SGD} & \texttt{Adam} \\
\midrule
is & the update you derived last week & a per-parameter adaptive step \\
typical \texttt{lr} & $10^{-2}$ to $10^{-1}$ & $10^{-3}$ \\
use when & you want to see the rule work & you want it to just train \\
\bottomrule
\end{tabular}
\end{center}

The learning rates are not interchangeable. Adam's default of $10^{-3}$ given to
SGD trains imperceptibly slowly; SGD's $10^{-1}$ given to Adam usually diverges.
When someone says a network ``did not learn'', the first question is which
optimizer, and the second is what learning rate.

\subsection{Reproducibility}

\begin{lstlisting}
torch.manual_seed(0)      # weight init, shuffling, dropout
\end{lstlisting}

Same caveat as Recitation 3, and it is worth repeating in this new setting:
seeding makes a run \emph{reproducible}, not \emph{meaningful}. A claim about an
architecture needs several seeds, because weight initialization is random and
the loss surface is not a bowl.

\section{Using a pretrained network}

You will spend more of this course \emph{reading} networks than training them.

\begin{lstlisting}
import torchvision
net = torchvision.models.alexnet(weights='DEFAULT')
net.eval()                                   # not training: turn off train mode

W = net.features[0].weight.detach().numpy()  # (64, 3, 11, 11)
\end{lstlisting}

\texttt{.detach()} removes the tensor from the autograd graph so it can become a
NumPy array; without it you get an error telling you to call it. The shape reads
as: 64 filters, each with 3 color channels, each $11\times11$ pixels. To display
one, move the channel dimension last and rescale to $[0,1]$:

\begin{lstlisting}
f = W[10].transpose(1, 2, 0)
plt.imshow((f - f.min()) / np.ptp(f))
\end{lstlisting}

Pretrained vision models expect their inputs normalized the way they were
trained --- a specific per-channel mean and standard deviation. Skip that step
and the filters still respond, but the numbers are not comparable to anything
published. Whenever you feed a pretrained model your own images, find the
normalization it was trained with and apply it.

\section{When it does not learn}

Work down this list in order. It is ordered by how often each one is the answer.

\begin{enumerate}
  \item \textbf{Is the loss going down at all?} Print it every epoch, on a log
        axis. ``Flat'' and ``slow'' look identical on a linear axis.
  \item \textbf{Did you call \texttt{opt.zero\_grad()}?} Symptom: the loss falls
        for a few epochs, then explodes.
  \item \textbf{Learning rate.} Try $\times 10$ and $\div 10$ before anything
        else. This fixes more cases than every other item combined.
  \item \textbf{Are the targets the right dtype and shape?} Class indices as
        \texttt{long}, shape $(N,)$ --- not one-hot, not $(N,1)$.
  \item \textbf{Is there a nonlinearity between the linear layers?} Without one
        the model is linear regardless of depth, and will plateau exactly where
        a linear model plateaus.
  \item \textbf{Did you apply softmax or ReLU to the logits?} Silent, and it
        halves your accuracy.
  \item \textbf{Can it overfit 10 examples?} Train on a tiny subset and try to
        drive the loss to zero. If it cannot, the bug is in the model or the
        loss. If it can, the model works and the problem is the data or the
        amount of training.
\end{enumerate}

\begin{keybox}
\ask{Item 7 is the one to remember.} ``Can it memorize a handful of examples?''
separates a broken pipeline from a hard problem in about thirty seconds, and it
is the first thing an experienced person does. A model that cannot overfit ten
points will never generalize from ten thousand.
\end{keybox}

\section{Exercises}

\ask{1.} You write \texttt{model(X)} and get \texttt{expected scalar type Float
but found Double}. What did you do, and what is the fix?

\ask{2.} Your loss decreases for 5 epochs, then grows without bound and reaches
\texttt{nan}. Give the two most likely causes and how to tell them apart.

\ask{3.} A classmate's classifier ends with
\texttt{nn.Linear(16, 10)} followed by \texttt{nn.ReLU()}, trained with
\texttt{nn.CrossEntropyLoss}. It reaches 60\% where yours reaches 95\%. Explain
what the ReLU did.

\ask{4.} \texttt{nn.Sequential(nn.Linear(64,16), nn.Linear(16,10))}. How many
parameters, and what function class can this model express?

\ask{5.} You extract features from a pretrained network inside a
\texttt{for} loop over 5,000 images and run out of memory. One line fixes it.
Which?

\ask{6.} Your model gets 100\% training accuracy on 10 examples and 12\% test
accuracy on 10,000. Is the pipeline broken?

\subsection*{Answers}

\ask{1.} You passed a NumPy \texttt{float64} array (or a tensor built from one)
into a layer expecting \texttt{float32}. Fix:
\texttt{torch.tensor(X, dtype=torch.float32)}.

\ask{2.} (i) A missing \texttt{opt.zero\_grad()}, so gradients accumulate and
each step grows; (ii) a learning rate too large. Tell them apart by fixing
\texttt{zero\_grad} first --- if it still diverges, the learning rate is the
cause. A missing \texttt{zero\_grad} typically shows a few good epochs first;
too large a learning rate usually misbehaves from the start.

\ask{3.} \texttt{CrossEntropyLoss} expects raw logits. ReLU clamps every
negative logit to zero, so the model can no longer express ``this class is
unlikely'' --- all the evidence \emph{against} a class is destroyed, leaving only
evidence for. Nothing errors.

\ask{4.} $(64{\times}16 + 16) + (16{\times}10 + 10) = 1{,}040 + 170 = 1{,}210$
parameters. With no nonlinearity between them it can express exactly the linear
maps from $\mathbb{R}^{64}$ to $\mathbb{R}^{10}$ --- the same class as a single
\texttt{nn.Linear(64, 10)}, with more parameters and no more power.

\ask{5.} Wrap the loop in \texttt{with torch.no\_grad():}. Without it, PyTorch
keeps the graph for all 5,000 forward passes.

\ask{6.} No --- that is the expected outcome of the overfit-ten-examples check,
and it tells you the model and loss are wired correctly. It says nothing about
generalization, which is a separate question and the subject of next week.

\section*{Cheat sheet}
\addcontentsline{toc}{section}{Cheat sheet}

\begin{center}
\begin{tabular}{@{}p{0.44\textwidth}p{0.5\textwidth}@{}}
\toprule
\textbf{You want} & \textbf{You write} \\
\midrule
NumPy $\rightarrow$ tensor             & \texttt{torch.tensor(a, dtype=torch.float32)} \\
tensor $\rightarrow$ NumPy             & \texttt{t.detach().numpy()} \\
add a batch dimension                  & \texttt{x[None]} \\
image for \texttt{imshow}              & \texttt{img.permute(1, 2, 0)} \\
a small MLP                            & \texttt{nn.Sequential(nn.Linear(..), nn.ReLU(), ..)} \\
count parameters                       & \texttt{sum(p.numel() for p in m.parameters())} \\
one training step                      & \texttt{zero\_grad; forward; loss; backward; step} \\
evaluate without recording             & \texttt{with torch.no\_grad():} \\
predicted class                        & \texttt{out.argmax(1)} \\
a Python float from a tensor           & \texttt{loss.item()} \\
a pretrained network                   & \texttt{torchvision.models.alexnet(weights='DEFAULT')} \\
same run twice                         & \texttt{torch.manual\_seed(0)} \\
\bottomrule
\end{tabular}
\end{center}

\end{document}
